一份读书笔记里有三条记录，分别写在相隔几百页的地方。第一条：梯度下降的步长上界由曲率的上界决定。第二条：Newton 法在拐点附近会发散。第三条：扩展 Kalman 滤波在急转弯处给出的协方差过分乐观。三条都记对了，三条都能复述，三条从来没有被放在一起过。

放在一起，它们是同一句话的三次记录：一阶近似只在一个有限的邻域里可信，而邻域的宽度由二阶项决定。知道这一点之后，三条记录塌缩成一条，而且多出一件原来没有的东西——下次在一个完全陌生的方法里看到``在当前点线性化''，不必等它出错就知道该问什么。

前面十六个部分是纵着切的：一个部分一个学科，各自把自己的对象、工具和边界讲完。这样切有它的道理，学科的内部逻辑是连贯的。代价是同一个想法在不同部分里换了名字，读者要在十六处分别记住它，而且很可能不知道那是同一处。这一单元横着切一刀：挑出五个真正在起作用的数学动作，把它们在各学科里的化身摆在同一页上，并说清每一次成立所依赖的条件。

\begin{center}
\begin{tikzpicture}[x=1cm,y=1cm,font=\scriptsize,>=stealth]
  % 列标题
  \node[font=\footnotesize,black!70] at (3.4,5.35) {微积分};
  \node[font=\footnotesize,black!70] at (6.0,5.35) {线性代数};
  \node[font=\footnotesize,black!70] at (8.6,5.35) {概率与统计};
  \node[font=\footnotesize,black!70] at (11.2,5.35) {优化与学习};
  % 竖列虚线
  \foreach \x in {3.4,6.0,8.6,11.2} \draw[black!18] (\x,0.2) -- (\x,5.05);
  % 五条横线
  \foreach \y/\name/\col in {4.6/线性化/blue,3.6/换基/teal,2.6/迭代/violet,1.6/不变量/olive,0.6/换序/orange}
    {
      \draw[\col!55!black,thick] (2.4,\y) -- (12.45,\y);
      \node[anchor=east,font=\footnotesize,\col!45!black] at (2.25,\y) {\name};
    }
  % 交点
  \foreach \x/\y/\t in {%
    3.4/4.6/切线, 6.0/4.6/Jacobian, 8.6/4.6/{delta 法}, 11.2/4.6/梯度下降,
    3.4/3.6/{Fourier 基}, 6.0/3.6/特征分解, 8.6/3.6/主成分, 11.2/3.6/低秩近似,
    3.4/2.6/Picard, 6.0/2.6/幂法, 8.6/2.6/MCMC, 11.2/2.6/值迭代,
    3.4/1.6/守恒量, 6.0/1.6/行列式, 8.6/1.6/鞅, 11.2/1.6/Lyapunov,
    3.4/0.6/控制收敛, 6.0/0.6/混合偏导, 8.6/0.6/Fubini, 11.2/0.6/重参数化}
    {\node[fill=white,inner sep=1.6pt,text=black!80] at (\x,\y) {\t};}
\end{tikzpicture}

\emph{每一条横线是一个数学动作，每一个交点是它在那个学科里用的名字。这张表要读者记住的不是二十个词，是五条横线——同一行上的四个名字，失效条件是同一个。}
\end{center}

\subsection{一条共同的句式}

五个动作都遵循同一个句式：\emph{先把问题换成一个便宜的形式，再交代这次换算在什么范围内有效。}前半句是它们被使用的理由，后半句是它们会出事的地方，而绝大多数事故发生在只做了前半句的时候。

线性化把非线性换成线性，代价是只在一个邻域内成立，邻域的宽度由曲率与斜率之比给出。换基把纠缠的线性作用换成独立的缩放，代价是那组基必须真的存在、而且被稳定地确定下来。迭代把解不出来的解换成一列近似的极限，代价是要说清在什么范数下压缩、以及那个不动点是不是你要的解。不变量把``检查所有可能''换成``检查一个数''，代价是找到合适的不变量没有通用方法。换序把两个操作的顺序调过来，代价是要防住质量逃逸。

这五句话的形状一样，而它们各自的``代价''那一半，正是前面十六个部分反复在讲的东西。这一单元不引入任何新结果，它只是把这些代价按动作而不是按学科重新归了一次档。

\subsection{五篇怎么安排}

按顺序读有一点好处，因为后一篇会用到前一篇的结论，但每一篇都能单独读。

\hyperref[sec:17-Synthesis/02-Recurring-Ideas/01]{``线性化：把弯的暂时当成直的''}处理出现次数最多的那个动作，从导数本身一直排到神经切核。它的核心是一个可以直接算的量——可信邻域的半径——以及八种表面上毫不相干的失效为什么都是这个半径被跨越。做优化、滤波、数值计算的读者，这一篇的回报最直接。

\hyperref[sec:17-Synthesis/02-Recurring-Ideas/02]{``换基：让复杂的作用变成独立的缩放''}回答两个问题：好基什么时候存在，以及它为什么就是好的。SVD 之所以普适，是因为它放弃了输入输出共用一组基；这一句想清楚，一批关于分解的记忆就可以不必背了。篇末处理一个高频误解——主成分是方差最大的方向，不是最有用的方向。

\hyperref[sec:17-Synthesis/02-Recurring-Ideas/03]{``迭代：把解写成一列近似的极限''}区分三种强度完全不同的收敛保证：压缩、单调、随机逼近。它最要紧的一节是最后一节，讲收敛到不动点为什么不等于收敛到你想要的东西——这一条在函数逼近与强化学习里造成过大量看起来像 bug 的现象。

\hyperref[sec:17-Synthesis/02-Recurring-Ideas/04]{``不变量：证明某件事做不到''}换了个方向：不去构造解，而去排除可能。它是这五篇里唯一一个主要产出否定性结论的动作，也是唯一一个必须承认``没有通用找法''的动作。

\hyperref[sec:17-Synthesis/02-Recurring-Ideas/05]{``换序：两个操作能不能交换顺序''}是全书正文的最后一篇。它把极限与积分、导数与期望、期望与最大这些日常操作背后的同一个条件挖出来，并给出一个框架永远不会报错的反例——支撑随参数移动。

\subsection{怎样使用这一单元}

不建议把这五篇当成需要记住的内容。它们的作用在识别：读一份陌生材料时，认出正在发生的是哪个动作，接下来该问什么就是确定的了。如果只带走一件事，带走这个提问的习惯——\emph{这个动作，在什么条件下合法。}

\hyperref[ch:139]{``按洞见查找：同一种结构怎样跨越学科''}是这一单元的索引式补充，按具体洞见给出跨部分的落点；\hyperref[ch:106]{``本质结构：对称性、低维性、信息几何与因果''}从另一个角度做同样的横切，它关心的是数据与模型里反复出现的\emph{结构}，而这一单元关心的是反复出现的\emph{动作}。两者可以对照着读。
